 \documentclass[12pt,a4paper]{article}
 \usepackage[utf8]{inputenc}
 \usepackage[T1]{fontenc}
 \usepackage[spanish]{babel}
 \usepackage{amsmath}
 \usepackage{amssymb}
 \usepackage{geometry}
 \usepackage{hyperref}
 \usepackage{bookmark}
 \usepackage{fancyhdr}
 \usepackage{enumitem}
 \usepackage{booktabs}
 \usepackage{titlesec}
 \geometry{margin=2.5cm}

 % Encabezado estilo notas de clase
 \setlength{\headheight}{15pt}
 \pagestyle{fancy}
 \fancyhf{}
 \fancyhead[L]{Microestructura, Eficiencia, Renta Fija y Derivados}
 \fancyhead[R]{Problem Set \#4}
 \fancyfoot[C]{\thepage}
 \fancyfoot[L]{TA Session \#4}
 \fancyfoot[R]{Jesus Villota Miranda}
 \renewcommand{\footrulewidth}{0.4pt}

 % Encabezado
 \title{\textbf{Problem Set 4}\\\large Microestructura, Eficiencia, Renta Fija y Derivados}
 \author{Jesus Villota Miranda\\CEMFI}
 \date{\today}

 \newcommand{\E}{\mathbb{E}}
 \newcommand{\V}{\mathbb{V}}
 \newcommand{\Cov}{\mathbb{C}\text{ov}}
 \newcommand{\R}{\mathbb{R}}
 \newcommand{\EUR}[1]{\ensuremath{#1~\text{EUR}}}

 % Formato de ejercicios
 \newenvironment{ejercicio}[1]{\subsection*{Ejercicio #1}}{}
 \newenvironment{solucion}{\paragraph{Solución.}}{}

 \begin{document}

 % \maketitle

 \section{Microestructura y negociación}

 \begin{ejercicio}{1 (Compra al margen y llamada de mantenimiento)}
 Una inversora compra \(400\) acciones de una compañía a \EUR{35} cada una financiando \EUR{4{,}500} con su bróker. El préstamo devenga un interés anual del \(9\%\).
 \begin{enumerate}[label=\alph*)]
   \item Calcule el margen inicial en euros y como porcentaje del valor de la posición.
   \item Si al cabo de un año el precio cae a \EUR{28} por acción, determine el valor del margen tras pagar intereses.
   \item Con un margen de mantenimiento del \(30\%\), ¿hay llamada de margen? Justifique con el porcentaje de margen.
   \item Calcule la rentabilidad de la inversora sobre su capital propio y compárela con la variación porcentual del precio.
 \end{enumerate}
 \end{ejercicio}
 \begin{solucion}
 La compra apalancada expone la rentabilidad del capital a los cambios de precio. El préstamo añade el coste financiero.
 \begin{enumerate}[label=\alph*)]
   \item El coste total es \(400\times 35 = \EUR{14{,}000}\). El margen inicial es \(14{,}000-4{,}500 = \EUR{9{,}500}\), que supone \(9{,}500/14{,}000 = 67{,}9\%\).
   \item Tras un año el valor de mercado es \(400\times 28 = \EUR{11{,}200}\). La deuda crece a \(4{,}500(1.09)=\EUR{4{,}905}\). El margen remanente es \(11{,}200-4{,}905 = \EUR{6{,}295}\).
   \item El porcentaje de margen es \(6{,}295/11{,}200 = 56{,}2\%\), por encima del \(30\%\); no hay llamada de margen.
   \item La rentabilidad sobre el capital propio es \((6{,}295-9{,}500)/9{,}500 = -0.337 = -33{,}7\%\). El precio cayó un \( (28-35)/35 = -20\% \). La pérdida se amplifica por el apalancamiento y los intereses.
 \end{enumerate}
 \end{solucion}

 \begin{ejercicio}{2 (Venta corta con dividendo y mantenimiento)}
 Un inversor abre una posición corta de \(800\) acciones a \EUR{48} con margen inicial del \(50\%\). La acción paga un dividendo de \EUR{1{,}50} dentro del año. Al cierre del año el precio sube a \EUR{55}.
 \begin{enumerate}[label=\alph*)]
   \item Calcule el efectivo obtenido y el margen inicial en euros.
   \item Determine el margen restante después del dividendo y del cambio de precio.
   \item Si el margen de mantenimiento es \(35\%\), indique si se produce llamada de margen.
   \item Compute la rentabilidad de la posición corta tomando el margen inicial como inversión.
 \end{enumerate}
 \end{ejercicio}
 \begin{solucion}
 En una venta corta el inversor recibe el efectivo de la venta pero debe aportar margen y reembolsar dividendos al prestamista de las acciones.
 \begin{enumerate}[label=\alph*)]
   \item El efectivo de la venta corta es \(800\times 48 = \EUR{38{,}400}\). El margen inicial del \(50\%\) es \(\EUR{19{,}200}\).
   \item La subida de precio genera una pérdida no realizada de \( (55-48)\times 800 = \EUR{5{,}600}\). El dividendo a devolver es \(1{,}50\times 800 = \EUR{1{,}200}\). El margen se reduce a \(19{,}200 - 5{,}600 - 1{,}200 = \EUR{12{,}400}\).
   \item El valor de la posición corta (pasivo) es \(55\times 800 = \EUR{44{,}000}\). El porcentaje de margen es \(12{,}400/44{,}000 = 28{,}2\%\), inferior al \(35\%\); hay llamada de margen.
   \item La rentabilidad es \((12{,}400-19{,}200)/19{,}200 = -35{,}4\%\). El precio subió un \(14{,}6\%\); la pérdida se amplifica por la posición corta y el dividendo pagado.
 \end{enumerate}
 \end{solucion}

 \begin{ejercicio}{3 (Libro de órdenes y costes de transacción)}
 Considere el siguiente libro de órdenes limitadas para la acción Soltec. La última operación fue a \EUR{62}.
 \begin{center}
 \begin{tabular}{cc|cc}
 \toprule
 \multicolumn{2}{c|}{Órdenes de compra} & \multicolumn{2}{c}{Órdenes de venta} \\
 Precio & Cantidad & Precio & Cantidad\\
 \midrule
 61{,}80 & 400 & 62{,}30 & 200\\
 61{,}50 & 600 & 63{,}10 & 300\\
 61{,}20 & 500 & 64{,}00 & 400\\
 60{,}80 & 900 & 65{,}50 & 600\\
 \bottomrule
 \end{tabular}
 \end{center}
 \begin{enumerate}[label=\alph*)]
   \item Determine el \emph{bid-ask spread} y el tamaño medio a ambos lados en el mejor nivel.
   \item Una orden de compra a mercado por \(300\) acciones, ¿a qué precio medio efectivo se ejecuta?
   \item ¿Cuál sería el precio medio y el coste total si la orden de compra fuera de \(900\) acciones?
   \item Calcule el coste de transacción implícito (en euros y como \% del valor nocional) de la orden de \(900\) acciones usando el último precio de \EUR{62} como referencia.
 \end{enumerate}
 \end{ejercicio}
 \begin{solucion}
 El coste de transacción incluye el spread y el impacto por consumir liquidez a precios peores.
 \begin{enumerate}[label=\alph*)]
   \item Mejor bid \(=61{,}80\) (400 acc.) y mejor ask \(=62{,}30\) (200 acc.). Spread: \(62{,}30-61{,}80 = \EUR{0{,}50}\). Tamaño medio en el mejor nivel: \((400+200)/2 = 300\) acciones.
   \item La orden de 300 compra primero 200 a \(62{,}30\) y 100 a \(63{,}10\). Precio medio ponderado: \((200\times 62{,}30 + 100\times 63{,}10)/300 = \EUR{62{,}57}\).
  \item Para 900 acciones: consume 200 a \(62{,}30\), 300 a \(63{,}10\) (queda 400), 400 a \(64{,}00\). Precio medio: \(\frac{200(62{,}30)+300(63{,}10)+400(64{,}00)}{900} = \EUR{63{,}32}\). Coste total pagado: \(900\times 63{,}32 = \EUR{56{,}990}\).
  \item Valor a precio de referencia \(62\): \(900\times 62 = \EUR{55{,}800}\). Coste implícito: \(56{,}990-55{,}800 = \EUR{1{,}190}\), equivalente a \(1{,}190/55{,}800 = 2{,}13\%\) del valor nocional.
 \end{enumerate}
 \end{solucion}

 \section{Eficiencia de mercado y eventos}

 \begin{ejercicio}{4 (Rentabilidad anormal y CAR con CAPM)}
 Una acción tiene \(\alpha=0\) y \(\beta=1{,}2\) estimados de forma estable. Durante tres meses consecutivos los rendimientos mensuales del mercado fueron \(3\%\), \(1\%\) y \(-2\%\). Los rendimientos observados de la acción fueron \(4{,}6\%\), \(2{,}0\%\) y \(-0{,}5\%\). Suponga \(r_f=0\).
 \begin{enumerate}[label=\alph*)]
   \item Calcule el rendimiento esperado en cada mes según el CAPM.
   \item Obtenga los rendimientos anormales mensuales.
   \item Compute el CAR (cumulative abnormal return) del trimestre y evalúe si indica noticia positiva o negativa.
 \end{enumerate}
 \end{ejercicio}
\begin{solucion}
El modelo CAPM establece que el rendimiento esperado de un activo es:
\[
E(r_i) = r_f + \beta_i [E(r_M) - r_f]
\]
Dado que \(\alpha=0\) y \(r_f=0\), la fórmula se simplifica a:
\[
E(r_i) = \beta_i \times r_M = 1{,}2 \times r_M
\]
donde \(r_M\) es el rendimiento del mercado observado en cada mes.

\begin{enumerate}[label=\alph*)]
  \item \textbf{Rendimientos esperados según CAPM:}
  
  Para cada mes, aplicamos la fórmula \(E(r) = 1{,}2 \times r_M\):
  \begin{itemize}
    \item \textbf{Mes 1:} \(E(r_1) = 1{,}2 \times 3\% = 3{,}6\%\)
    \item \textbf{Mes 2:} \(E(r_2) = 1{,}2 \times 1\% = 1{,}2\%\)
    \item \textbf{Mes 3:} \(E(r_3) = 1{,}2 \times (-2\%) = -2{,}4\%\)
  \end{itemize}
  
  \item \textbf{Rendimientos anormales mensuales:}
  
  El rendimiento anormal (abnormal return) se define como la diferencia entre el rendimiento observado y el rendimiento esperado según el modelo:
  \[
  AR_t = r_{t,\text{observado}} - E(r_t)
  \]
  
  Calculamos para cada mes:
  \begin{itemize}
    \item \textbf{Mes 1:} \(AR_1 = 4{,}6\% - 3{,}6\% = 1{,}0\%\)
    \item \textbf{Mes 2:} \(AR_2 = 2{,}0\% - 1{,}2\% = 0{,}8\%\)
    \item \textbf{Mes 3:} \(AR_3 = -0{,}5\% - (-2{,}4\%) = -0{,}5\% + 2{,}4\% = 1{,}9\%\)
  \end{itemize}
  
  \item \textbf{CAR (Cumulative Abnormal Return) del trimestre:}
  
  El CAR acumula los rendimientos anormales a lo largo de la ventana de evento:
  \[
  CAR = \sum_{t=1}^{3} AR_t = AR_1 + AR_2 + AR_3
  \]
  
  Sustituyendo los valores calculados:
  \[
  CAR = 1{,}0\% + 0{,}8\% + 1{,}9\% = 3{,}7\%
  \]
  
  \textbf{Interpretación:} Un CAR positivo de \(3{,}7\%\) indica que la acción ha tenido rendimientos sistemáticamente superiores a los esperados según el CAPM durante el trimestre. Esto sugiere que:
  \begin{itemize}
    \item Puede haber habido noticias favorables específicas de la compañía no capturadas por el modelo de mercado.
    \item La acción podría haber estado infravalorada al inicio de la ventana de observación.
    \item Puede existir un alpha positivo no modelado (aunque se asumió \(\alpha=0\)).
  \end{itemize}
\end{enumerate}
\end{solucion}

 \section{Renta fija}

 \begin{ejercicio}{6 (Comparación de bonos: cupón cero vs cupón anual)}
 Considere dos bonos con nominal \EUR{1{,}000} y vencimiento a 2 años:
 \begin{itemize}
   \item \textbf{Bono A:} Bono cupón cero (zero coupon bond)
   \item \textbf{Bono B:} Bono con cupón anual del \(5\%\)
 \end{itemize}
 Los tipos spot anuales son: \(r_1=3{,}0\%\) y \(r_2=3{,}5\%\). Los bonos fueron adquiridos a los siguientes precios de compra: Bono A a \EUR{925} y Bono B a \EUR{1{,}035}.
 \begin{enumerate}[label=\alph*)]
   \item Calcule el precio teórico de ambos bonos usando la curva spot.
  \item Asumiendo que los bonos se mantienen hasta el vencimiento, calcule la TIR (yield to maturity) de cada bono. La TIR es la tasa única de descuento que iguala el precio de compra con el valor presente de los flujos futuros. Para el Bono B, resuelva la ecuación cuadrática correspondiente para calcular la TIR. Comente la relación entre la TIR y los tipos spot, y explique por qué la TIR del Bono A difiere del tipo spot \(r_2\) cuando el precio de compra no coincide con el teórico.
   \item Calcule la duración de Macaulay de ambos bonos usando sus respectivas TIR como tasa de descuento uniforme.
   \item La nueva curva de tipos muestra que \(r_1 = 3{,}5\%\) y el tipo forward implícito \(f_{1,1} = 4{,}5\%\). Calcule primero el tipo spot \(r_2\) a partir de \(r_1\) y \(f_{1,1}\). Luego, calcule el nuevo precio teórico de ambos bonos usando la nueva curva spot.
 \end{enumerate}
 \end{ejercicio}
 \begin{solucion}
 Este ejercicio ilustra cómo la estructura de pagos afecta la sensibilidad del precio a cambios en los tipos de interés.
 \begin{enumerate}[label=\alph*)]
   \item \textbf{Precio del Bono A (cupón cero):}
   
   El bono cupón cero solo paga el nominal al vencimiento. Su precio es:
   \[
   P_A = \frac{1{,}000}{(1+r_2)^2} = \frac{1{,}000}{(1{,}035)^2} = \frac{1{,}000}{1{,}071225} = \EUR{933{,}51}
   \]
   
   \textbf{Precio del Bono B (con cupón anual del 5\%):}
   
   El bono paga un cupón de \EUR{50} al final del año 1 y \EUR{1{,}050} (cupón + nominal) al final del año 2:
   \[
   P_B = \frac{50}{1+r_1} + \frac{1{,}050}{(1+r_2)^2} = \frac{50}{1{,}03} + \frac{1{,}050}{(1{,}035)^2} = 48{,}54 + 980{,}18 = \EUR{1{,}028{,}72}
   \]
   
   \item \textbf{TIR del Bono A:}
   
   La TIR se calcula usando el precio de compra \(P_A^{\text{compra}} = \EUR{925}\):
   \[
   925 = \frac{1{,}000}{(1+\text{TIR}_A)^2}
   \]
   
   Despejando: \((1+\text{TIR}_A)^2 = \frac{1{,}000}{925} = 1{,}081081\), por lo que \(1+\text{TIR}_A = \sqrt{1{,}081081} = 1{,}03976\).
   
   Por tanto, \(\text{TIR}_A = 3{,}976\%\). 
   
   \textbf{Comentario:} La TIR (3{,}976\%) es mayor que el tipo spot \(r_2 = 3{,}5\%\) porque el bono se compró a un precio inferior al teórico (\EUR{925} < \EUR{933{,}51}). Si el precio de compra hubiera coincidido con el precio teórico calculado con la curva spot, la TIR del bono cupón cero habría sido exactamente igual a \(r_2\), ya que el único pago se produce al vencimiento y el precio teórico se calcula descontando ese pago precisamente con \(r_2\). La diferencia entre el precio de mercado y el precio teórico refleja factores como las condiciones de mercado, la liquidez, o diferencias en las expectativas de los inversores respecto a la curva spot de referencia.
   
   \textbf{TIR del Bono B:}
   
   La TIR se calcula usando el precio de compra \(P_B^{\text{compra}} = \EUR{1{,}035}\):
   \[
   1{,}035 = \frac{50}{1+\text{TIR}_B} + \frac{1{,}050}{(1+\text{TIR}_B)^2}
   \]
   
   Multiplicando ambos lados por \((1+\text{TIR}_B)^2\):
   \[
   1{,}035(1+\text{TIR}_B)^2 = 50(1+\text{TIR}_B) + 1{,}050
   \]
   
   Haciendo el cambio de variable \(x = 1+\text{TIR}_B\):
   \[
   1{,}035x^2 - 50x - 1{,}050 = 0
   \]
   
   Resolviendo la ecuación cuadrática con \(a = 1{,}035\), \(b = -50\), \(c = -1{,}050\):
   \[
   \Delta = 50^2 + 4(1{,}035)(1{,}050) = 2{,}500 + 4{,}347{,}000 = 4{,}349{,}500
   \]
   \[
   \sqrt{\Delta} = \sqrt{4{,}349{,}500} = 2{,}085{,}55
   \]
   \[
   x = \frac{50 + 2{,}085{,}55}{2{,}070} = \frac{2{,}135{,}55}{2{,}070} = 1{,}0317
   \]
   
   (Descartamos la raíz negativa porque \(x = 1+\text{TIR}_B > 0\)).
   
   Por tanto, \(\text{TIR}_B = 1{,}0317 - 1 = 3{,}17\%\). 
   
   \textbf{Comentario:} La TIR (3{,}17\%) es menor que el cupón del \(5\%\) porque el precio de compra está por encima del nominal (\EUR{1{,}035} > \EUR{1{,}000}). Además, la TIR es menor que una tasa promedio ponderada de los tipos spot porque el precio de compra (\EUR{1{,}035}) es superior al precio teórico calculado con la curva spot (\EUR{1{,}028{,}72}), lo que reduce la rentabilidad efectiva para el inversor.
   
   \item \textbf{Duración de Macaulay:}
   
   Calculamos la duración usando las TIR obtenidas en el apartado anterior, que son las tasas de descuento que igualan los precios de compra con los flujos futuros.
   
   \textbf{Bono A:} Usando la TIR \(\text{TIR}_A = 3{,}976\%\) y el precio de compra \(P_A^{\text{compra}} = \EUR{925}\):
   \[
   PV(\text{pago año 2}) = \frac{1{,}000}{(1{,}03976)^2} = 925
   \]
   
   La duración de Macaulay es:
   \[
   D_A = \frac{2 \times 925}{925} = 2\ \text{años}
   \]
   
   Para un bono cupón cero, la duración siempre coincide con el tiempo hasta el vencimiento, independientemente de la TIR.
   
   \textbf{Bono B:} Usando la TIR \(\text{TIR}_B = 3{,}17\%\) y el precio de compra \(P_B^{\text{compra}} = \EUR{1{,}035}\):
   \begin{align*}
   PV(C_1) &= \frac{50}{1{,}0317} = 48{,}46\\
   PV(C_2+N) &= \frac{1{,}050}{(1{,}0317)^2} = \frac{1{,}050}{1{,}0644} = 986{,}54\\
   P_B &= 48{,}46 + 986{,}54 = 1{,}035
   \end{align*}
   
   La duración de Macaulay es:
   \[
   D_B = \frac{1 \times 48{,}46 + 2 \times 986{,}54}{1{,}035} = \frac{2{,}021{,}54}{1{,}035} = 1{,}953\ \text{años}
   \]
   
   La duración del bono con cupón es menor que su vencimiento porque recibe pagos antes del vencimiento.
   
   \item \textbf{Cálculo de \(r_2\) a partir de \(r_1\) y \(f_{1,1}\):}
   
   La relación entre tipos spot y forward viene dada por la condición de no arbitraje:
   \[
   (1+r_2)^2 = (1+r_1)(1+f_{1,1})
   \]
   
   Despejando \(r_2\):
   \[
   (1+r_2)^2 = (1+0{,}035)(1+0{,}045) = 1{,}035 \times 1{,}045 = 1{,}081575
   \]
   
   Por tanto:
   \[
   1+r_2 = \sqrt{1{,}081575} = 1{,}0400
   \]
   
   Así, \(r_2 = 4{,}0\%\).
   
   \textbf{Cálculo de los nuevos precios teóricos:}
   
   Ahora calculamos los nuevos precios teóricos usando la nueva curva spot (\(r_1 = 3{,}5\%\), \(r_2 = 4{,}0\%\)):
   \begin{align*}
   P_A^{\text{nuevo}} &= \frac{1{,}000}{(1+r_2)^2} = \frac{1{,}000}{(1{,}04)^2} = \frac{1{,}000}{1{,}0816} = \EUR{924{,}56}\\
   P_B^{\text{nuevo}} &= \frac{50}{1+r_1} + \frac{1{,}050}{(1+r_2)^2} = \frac{50}{1{,}035} + \frac{1{,}050}{(1{,}04)^2} = 48{,}31 + 970{,}78 = \EUR{1{,}019{,}09}
   \end{align*}
   
   \textbf{Observación:} Ambos precios teóricos han disminuido respecto a los precios teóricos originales (\(P_A = \EUR{933{,}51}\), \(P_B = \EUR{1{,}028{,}72}\)), lo cual es consistente con el aumento en los tipos de interés. El bono cupón cero muestra una mayor variación absoluta en precio porque todos sus flujos se reciben al vencimiento, mientras que el bono con cupón recibe pagos intermedios que se descuentan a un tipo menor (\(r_1 = 3{,}5\%\)).
 \end{enumerate}
 \end{solucion}

 \section{Derivados: forwards}

 \begin{ejercicio}{8 (Curva spot y tipos forward implícitos)}
 Se observa la siguiente curva de tipos spot anuales (compuestos anualmente):
 \begin{center}
 \begin{tabular}{cc}
 \toprule
 Vencimiento (años) & Tipo spot anual (\%) \\
 \midrule
 1 & \(3{,}0\) \\
 2 & \(3{,}5\) \\
 3 & \(4{,}0\) \\
 4 & \(4{,}2\) \\
 5 & \(4{,}3\) \\
 \bottomrule
 \end{tabular}
 \end{center}
 \begin{enumerate}[label=\alph*)]
   \item Calcule el tipo forward implícito \(f_{1,1}\) (forward de 1 año que comienza dentro de 1 año).
   \item Calcule el tipo forward implícito \(f_{2,1}\) (forward de 1 año que comienza dentro de 2 años).
   \item Calcule el tipo forward implícito \(f_{1,2}\) (forward de 2 años que comienza dentro de 1 año).
   \item Si un inversor quiere asegurar hoy un tipo de interés del \(4{,}5\%\) para un préstamo de 1 año que comenzará dentro de 2 años, ¿le conviene más contratar un forward o esperar? Justifique su respuesta comparando con el tipo forward implícito.
   \item Explique intuitivamente por qué el tipo forward \(f_{1,1}\) es mayor que el tipo spot a 1 año \(r_1\) cuando la curva tiene pendiente positiva.
 \end{enumerate}
 \end{ejercicio}
 \begin{solucion}
 Los tipos forward implícitos se derivan de la curva spot usando la condición de no arbitraje: invertir a largo plazo debe ser equivalente a invertir a corto y luego reinvertir al tipo forward.
 \begin{enumerate}[label=\alph*)]
   \item El tipo forward \(f_{1,1}\) es el tipo de interés para el periodo entre \(t=1\) y \(t=2\). La condición de no arbitraje establece:
   \[
   (1+r_2)^2 = (1+r_1)(1+f_{1,1})
   \]
   Despejando:
   \[
   1+f_{1,1} = \frac{(1+r_2)^2}{1+r_1} = \frac{(1{,}035)^2}{1{,}03} = \frac{1{,}071225}{1{,}03} = 1{,}0400
   \]
   Por tanto, \(f_{1,1} = 4{,}00\%\).
   
   \item El tipo forward \(f_{2,1}\) es el tipo para el periodo entre \(t=2\) y \(t=3\):
   \[
   (1+r_3)^3 = (1+r_2)^2(1+f_{2,1})
   \]
   Despejando:
   \[
   1+f_{2,1} = \frac{(1+r_3)^3}{(1+r_2)^2} = \frac{(1{,}04)^3}{(1{,}035)^2} = \frac{1{,}124864}{1{,}071225} = 1{,}0500
   \]
   Por tanto, \(f_{2,1} = 5{,}00\%\).
   
   \item El tipo forward \(f_{1,2}\) es el tipo para el periodo entre \(t=1\) y \(t=3\) (2 años):
   \[
   (1+r_3)^3 = (1+r_1)(1+f_{1,2})^2
   \]
   Despejando:
   \[
   (1+f_{1,2})^2 = \frac{(1+r_3)^3}{1+r_1} = \frac{(1{,}04)^3}{1{,}03} = \frac{1{,}124864}{1{,}03} = 1{,}0921
   \]
   Por tanto, \(1+f_{1,2} = \sqrt{1{,}0921} = 1{,}0450\), así \(f_{1,2} = 4{,}50\%\).
   
   \textbf{Nota:} Este es el tipo forward anualizado para un periodo de 2 años. Si queremos el tipo efectivo para los 2 años completos, sería \((1{,}045)^2 - 1 = 9{,}20\%\), pero normalmente se reporta como tasa anualizada.
   
   \item El tipo forward implícito \(f_{2,1}\) calculado en el apartado b) es \(5{,}00\%\). Si el inversor puede asegurar un tipo del \(4{,}5\%\) mediante un contrato forward, esto es mejor que el tipo forward implícito de \(5{,}00\%\), ya que pagará menos intereses. 
   
   Alternativamente, si espera y toma el préstamo al tipo spot vigente en \(t=2\), pagará aproximadamente el tipo forward implícito (si las expectativas se cumplen), que es \(5{,}00\%\). Por tanto, le conviene contratar el forward al \(4{,}5\%\).
   
   \textbf{Interpretación:} El tipo forward implícito de \(5{,}00\%\) refleja las expectativas del mercado sobre el tipo spot futuro. Si el inversor puede asegurar un tipo menor (\(4{,}5\%\)), está obteniendo una mejor condición que la esperada por el mercado.
   
   \item Cuando la curva spot tiene pendiente positiva (tipos a largo plazo mayores que a corto plazo), el tipo forward \(f_{1,1}\) debe ser mayor que \(r_1\) para compensar la diferencia entre \(r_1\) y \(r_2\).
   
   \textbf{Intuición:} Si \(r_2 > r_1\), invertir durante 2 años directamente debe ofrecer la misma rentabilidad que invertir 1 año y luego reinvertir. Como el segundo año ``compensa'' la menor rentabilidad del primer año, el tipo forward del segundo año (\(f_{1,1}\)) debe ser mayor que \(r_1\) para que la rentabilidad total de ambas estrategias sea igual.
   
   Matemáticamente, de \((1+r_2)^2 = (1+r_1)(1+f_{1,1})\), si \(r_2 > r_1\), entonces \((1+r_2)^2 > (1+r_1)^2\), lo que implica que \((1+f_{1,1}) > (1+r_1)\), es decir, \(f_{1,1} > r_1\).
   
   \textbf{Ejemplo numérico:} Con \(r_1 = 3{,}0\%\) y \(r_2 = 3{,}5\%\), calculamos \(f_{1,1} = 4{,}0\%\), que efectivamente es mayor que \(r_1 = 3{,}0\%\). El tipo forward ``sube'' para compensar el tipo spot más bajo del primer año.
 \end{enumerate}
 \end{solucion}

\begin{ejercicio}{9 (Libro de órdenes y bonos corporativos AirBaltic)}
El bono corporativo de AirBaltic fue emitido el 15 de noviembre de 2024 con nominal de \EUR{1{,}000} cada uno, paga cupones trimestrales a una tasa anual del \(14{,}5\%\) y madura el 15 de noviembre de 2028. El bono es callable. El libro de órdenes muestra las siguientes posiciones (precios en \% del nominal, sin cupón corrido):
\begin{center}
\begin{tabular}{cc|cc}
\toprule
\multicolumn{2}{c|}{Órdenes de compra} & \multicolumn{2}{c}{Órdenes de venta} \\
Precio & Cantidad (nom.) & Precio & Cantidad (nom.)\\
\midrule
100{,}40 & 200{,}000 & 101{,}00 & 150{,}000\\
100{,}25 & 250{,}000 & 101{,}20 & 150{,}000\\
99{,}90  & 300{,}000 & 101{,}50 & 200{,}000\\
         &           & 102{,}00 & 100{,}000\\
\bottomrule
\end{tabular}
\end{center}
\begin{enumerate}[label=\alph*)]
  \item El 15 de diciembre de 2025, Jesus envía una orden de compra a mercado por nominal \EUR{350{,}000}. ¿A qué precio efectivo (precio medio ponderado en \% del nominal) adquiere el bono? Calcule también el cupón corrido que debe pagar en la fecha de compra y el desembolso total en euros.
  \item Si Jesus mantuviera el bono hasta su vencimiento (15 de noviembre de 2028), plantee la ecuación que debería satisfacer la TIR (YTM) anual del bono. Considere que el bono paga cupones trimestrales y devuelve el \(100\%\) del nominal al vencimiento, y que dado que el bono se compró el 15 de diciembre de 2025, los primeros pagos no coinciden con trimestres completos desde la compra.
  \item Sin embargo, Jesus decide vender el bono el 15 de abril de 2026. En esa fecha, el libro de órdenes muestra las siguientes posiciones (precios en \% del nominal, sin cupón corrido):
  \begin{center}
  \begin{tabular}{cc|cc}
  \toprule
  \multicolumn{2}{c|}{Órdenes de compra} & \multicolumn{2}{c}{Órdenes de venta} \\
  Precio & Cantidad (nom.) & Precio & Cantidad (nom.)\\
  \midrule
  99{,}20 & 100{,}000 & 99{,}80 & 200{,}000\\
  99{,}00 & 150{,}000 & 100{,}00 & 250{,}000\\
  98{,}75 & 200{,}000 & 100{,}30 & 300{,}000\\
  98{,}40 & 150{,}000 &         &          \\
  \bottomrule
  \end{tabular}
  \end{center}
  Calcule el precio efectivo de venta que obtiene Jesus al vender los \EUR{350{,}000} nominales mediante una orden de venta a mercado. Luego, determine la rentabilidad anual (TIR) obtenida por Jesus en esta inversión, considerando que al vender en abril tiene derecho a cobrar la parte proporcional del cupón trimestral (2/3 del cupón, ya que faltan 30 días para el pago del 15 de mayo).
\end{enumerate}
\end{ejercicio}
\begin{solucion}
Una orden a mercado barre los mejores asks hasta completar el tamaño; el cupón corrido se calcula proporcionalmente al tiempo transcurrido desde el último pago.
\begin{enumerate}[label=\alph*)]
  \item La orden consume \(150{,}000\) al \(101{,}00\%\), \(150{,}000\) al \(101{,}20\%\) y los \(50{,}000\) restantes al \(101{,}50\%\). El precio medio ponderado es
  \[
  \bar{P} = \frac{150{,}000(101{,}00)+150{,}000(101{,}20)+50{,}000(101{,}50)}{350{,}000} = 101{,}16\%\ \text{del nominal}.
  \]
  Cada bono (nominal \EUR{1{,}000}) cuesta \(1{,}011{,}6\) euros. El último cupón se pagó el 15 de noviembre de 2025, por lo que el cupón corrido es la parte proporcional del cupón trimestral correspondiente a 1 mes (1/3 del trimestre). Cupón trimestral = \EUR{36{,}25}, cupón corrido = \(36{,}25 \times 1/3 = \EUR{12{,}08}\). Desembolso total: \(350{,}000 \times (1{,}0116 + 0{,}01208) = \EUR{354{,}060} + \EUR{4{,}228} = \EUR{358{,}288}\).
  \item \textbf{Planteamiento de la TIR:}
  
  Primero identificamos el horizonte temporal. Desde la fecha de compra (15 de diciembre de 2025) hasta el vencimiento (15 de noviembre de 2028) transcurren exactamente \textbf{35 meses} (2 años y 11 meses), no 3 años.
  
  Los cupones trimestrales se pagan en fechas fijas del calendario: 15 de febrero, mayo, agosto y noviembre de cada año. Desde la compra, Jesus recibirá los siguientes pagos:
  
  \begin{itemize}
    \item \textbf{Cupón 1:} 15 feb 2026 (2 meses = 2/12 años después)
    \item \textbf{Cupón 2:} 15 may 2026 (5 meses = 5/12 años)
    \item \textbf{Cupón 3:} 15 ago 2026 (8 meses = 8/12 años)
    \item \textbf{Cupón 4:} 15 nov 2026 (11 meses = 11/12 años)
    \item \textbf{Cupón 5:} 15 feb 2027 (14 meses = 14/12 años)
    \item \textbf{Cupón 6:} 15 may 2027 (17 meses = 17/12 años)
    \item \textbf{Cupón 7:} 15 ago 2027 (20 meses = 20/12 años)
    \item \textbf{Cupón 8:} 15 nov 2027 (23 meses = 23/12 años)
    \item \textbf{Cupón 9:} 15 feb 2028 (26 meses = 26/12 años)
    \item \textbf{Cupón 10:} 15 may 2028 (29 meses = 29/12 años)
    \item \textbf{Cupón 11:} 15 ago 2028 (32 meses = 32/12 años)
    \item \textbf{Cupón 12 + Nominal:} 15 nov 2028 (35 meses = 35/12 años)
  \end{itemize}
  
  El desembolso inicial efectivo fue de \EUR{1{,}023{,}68} por bono (precio de \(101{,}16\%\) del nominal más cupón corrido de \EUR{12{,}08}). La TIR anual \(y\) es la tasa de descuento que iguala el valor presente de todos los flujos futuros con este desembolso inicial:
  
  \begin{align*}
  1{,}023{,}68 &= \frac{36{,}25}{(1+y)^{2/12}} + \frac{36{,}25}{(1+y)^{5/12}} + \frac{36{,}25}{(1+y)^{8/12}} + \frac{36{,}25}{(1+y)^{11/12}}\\
  &\quad + \frac{36{,}25}{(1+y)^{14/12}} + \frac{36{,}25}{(1+y)^{17/12}} + \frac{36{,}25}{(1+y)^{20/12}} + \frac{36{,}25}{(1+y)^{23/12}}\\
  &\quad + \frac{36{,}25}{(1+y)^{26/12}} + \frac{36{,}25}{(1+y)^{29/12}} + \frac{36{,}25}{(1+y)^{32/12}} + \frac{1{,}036{,}25}{(1+y)^{35/12}}
  \end{align*}
  
  De forma compacta:
  \[
  1{,}023{,}68 = \sum_{k=1}^{11} \frac{36{,}25}{(1+y)^{(3k-1)/12}} + \frac{1{,}036{,}25}{(1+y)^{35/12}}
  \]
  
  donde \((3k-1)\) genera la secuencia \(2, 5, 8, 11, \ldots, 32\) que representa los meses hasta cada cupón.
  
  \textbf{Nota práctica:} Esta ecuación no tiene solución cerrada y debe resolverse numéricamente (mediante Newton-Raphson o tanteo). El cupón anual nominal del \(14{,}5\%\) y el precio ligeramente superior al par sugieren que la TIR será cercana pero algo inferior al \(14{,}5\%\), ajustada por el momento exacto de cada pago.
  \item \textbf{Paso 1: Calcular el precio efectivo de venta}
  
  Una orden de venta a mercado consume los mejores bids (órdenes de compra) disponibles. Para vender \EUR{350{,}000} nominales, Jesus barre el libro:
  \begin{itemize}
    \item \EUR{100{,}000} al \(99{,}20\%\)
    \item \EUR{150{,}000} al \(99{,}00\%\)
    \item \EUR{100{,}000} restantes al \(98{,}75\%\)
  \end{itemize}
  
  El precio medio ponderado de venta es:
  \[
  \bar{P}_{\text{venta}} = \frac{100{,}000(99{,}20) + 150{,}000(99{,}00) + 100{,}000(98{,}75)}{350{,}000} = \frac{34{,}645{,}000}{350{,}000} = 99{,}01\% 
  \]
  Por cada bono (nominal \EUR{1{,}000}), Jesus recibe \EUR{990{,}14}.
  
  \textbf{Paso 2: Calcular el cupón corrido al vender}
  
  La venta se produce el 15 de abril de 2026. El último cupón se pagó el 15 de febrero de 2026, y el siguiente cupón se pagará el 15 de mayo de 2026. Han transcurrido exactamente 2 meses del trimestre (60 días de 90), por lo que Jesus tiene derecho a:
  \[
  \text{Cupón corrido} = 36{,}25 \times \frac{2}{3} = \EUR{24{,}17}
  \]
  
  \textbf{Paso 3: Identificar los flujos de caja}
  
  La cronología completa de la inversión es:
  \begin{itemize}
    \item \textbf{15 dic 2025} (\(t=0\)): Inversión inicial \(FC_0 = -\EUR{1{,}023{,}68}\) (precio \(101{,}16\%\) más cupón corrido \EUR{12{,}08})
    \item \textbf{15 feb 2026} (\(t=2/12\) años): Recibe cupón \(FC_1 = +\EUR{36{,}25}\)
    \item \textbf{15 abr 2026} (\(t=4/12\) años): Vende el bono por \(990{,}14 + 24{,}17 = \EUR{1{,}014{,}31}\), así \(FC_2 = +\EUR{1{,}014{,}31}\)
  \end{itemize}
  
  \textbf{Paso 4: Plantear la ecuación de valor presente}
  
  La TIR anual se obtiene igualando el valor presente de los flujos futuros con la inversión inicial:
  \[
  FC_0 = \frac{FC_1}{(1+TIR)^{2/12}} + \frac{FC_2}{(1+TIR)^{4/12}}
  \]
  
  Sustituyendo los valores:
  \[
  1{,}023{,}68 = \frac{36{,}25}{(1+TIR)^{2/12}} + \frac{1{,}014{,}31}{(1+TIR)^{4/12}}
  \]
  
  \textbf{Paso 5: Simplificar mediante cambio de variable}
  
  Definimos \(x := (1+TIR)^{2/12} = (1+TIR)^{1/6}\). Entonces:
  \begin{itemize}
    \item \((1+TIR)^{2/12} = x\)
    \item \((1+TIR)^{4/12} = x^2\)
  \end{itemize}
  
  La ecuación se transforma en:
  \[
  1{,}023{,}68 = \frac{36{,}25}{x} + \frac{1{,}014{,}17}{x^2}
  \]
  
  Multiplicando ambos lados por \(x^2\):
  \[
  1{,}023{,}68 x^2 = 36{,}25 x + 1{,}014{,}17
  \]
  
  Reordenando:
  \[
  1{,}023{,}68 x^2 - 36{,}25 x - 1{,}014{,}31 = 0
  \]
  
  \textbf{Paso 6: Resolver la ecuación cuadrática}
  
  Aplicamos la fórmula cuadrática \(x = \frac{-b \pm \sqrt{b^2 - 4ac}}{2a}\) con \(a = 1{,}023{,}68\), \(b = -36{,}25\), \(c = -1{,}014{,}17\):
  \[
  \Delta = (-36{,}25)^2 - 4(1{,}023{,}68)(-1{,}014{,}31) = 1{,}314{,}06 + 4{,}153{,}193{,}72 = 4{,}154{,}507{,}78
  \]
  \[
  \sqrt{\Delta} = 2{,}038{,}26
  \]
  \[
  x = \frac{36{,}25 + 2{,}038{,}26}{2(1{,}023{,}68)} = \frac{2{,}074{,}51}{2{,}047{,}36} = 1{,}0133
  \]
  
  (Descartamos la raíz negativa porque \(x\) debe ser positivo.)
  
  \textbf{Paso 7: Recuperar la TIR}
  
  Recordando que \(x = (1+TIR)^{1/6}\), tenemos:
  \[
  1 + TIR = x^6 = (1{,}0132)^6
  \]
  
  Calculando:
  \[
  (1{,}0133)^6 = 1{,}0824
  \]
  
  Por tanto:
  \[
  TIR = 1{,}0824 - 1 = 0{,}0824 = 8{,}24\%
  \]
  
  	extbf{Conclusión:} Jesus obtuvo una rentabilidad anualizada (TIR) del \textbf{8{,}24\%} en esta operación. Aunque el bono paga cupones nominales del \(14{,}5\%\) anual, la rentabilidad efectiva es menor porque: (i) compró el bono ligeramente por encima del par (\(101{,}16\%\)), (ii) lo vendió por debajo del par (\(99{,}01\%\)), sufriendo una pérdida de capital del \(2{,}15\%\) sobre el precio de compra, y (iii) solo recibió un cupón completo más el cupón corrido proporcional durante el periodo de tenencia de 4 meses. Esta rentabilidad del \(8{,}24\%\) sigue siendo atractiva en comparación con inversiones sin riesgo, reflejando el riesgo crediticio de AirBaltic.
\end{enumerate}
\end{solucion}

\section{Derivados: forwards, futures y opciones}

\begin{ejercicio}{10 (Valoración de forwards y payoffs de opciones)}
Una acción de la empresa TechCorp cotiza hoy a \EUR{45} y no paga dividendos. El tipo de interés libre de riesgo anual es del \(4\%\).
\begin{enumerate}[label=\alph*)]
  \item Calcule el precio forward teórico (precio de ejercicio \(K\)) para un contrato forward que vence dentro de 2 años usando el principio de no arbitraje.
  \item Si hace 6 meses se firmó un contrato forward para comprar la acción a \EUR{48} dentro de 2 años, y hoy la acción cotiza a \EUR{50}, calcule el valor actual del contrato forward para el comprador (posición larga).
  \item Un inversor compra una opción de compra (call) europea sobre TechCorp con precio de ejercicio \EUR{48} y vencimiento en 1 año, pagando una prima de \EUR{3{,}50}. Calcule el beneficio neto del inversor si al vencimiento el precio de la acción es: (i) \EUR{52}, (ii) \EUR{48}, (iii) \EUR{44}.
  \item El mismo inversor compra también una opción de venta (put) europea sobre TechCorp con precio de ejercicio \EUR{48} y vencimiento en 1 año, pagando una prima de \EUR{2{,}00}. Calcule el beneficio neto total si al vencimiento el precio de la acción es \EUR{44} (considerando ambas opciones).
\end{enumerate}
\end{ejercicio}
\begin{solucion}
Los derivados se valoran usando el principio de no arbitraje: una cartera réplica debe tener el mismo valor que el derivado.
\begin{enumerate}[label=\alph*)]
  \item \textbf{Precio forward teórico:}
  
  El precio forward teórico se calcula usando la cartera réplica: comprar la acción hoy y financiarla vendiendo en corto un bono cupón cero que paga \(K\) al vencimiento.
  
  El valor de la cartera réplica hoy es:
  \[
  F_0 = S_0 - K(1+r_f)^{-T}
  \]
  
  Para que el contrato forward tenga valor cero al inicio (condición de no arbitraje), debe cumplirse:
  \[
  0 = S_0 - K(1+r_f)^{-T}
  \]
  
  Despejando \(K\):
  \[
  K = S_0(1+r_f)^T = 45 \times (1{,}04)^2 = 45 \times 1{,}0816 = \EUR{48{,}67}
  \]
  
  \textbf{Interpretación:} El precio forward de \EUR{48{,}67} es el precio al que se puede comprar la acción dentro de 2 años sin riesgo de arbitraje, considerando el coste de oportunidad del dinero (tipo de interés del \(4\%\)).
  
  \item \textbf{Valor del contrato forward existente:}
  
  El contrato forward se firmó hace 6 meses con precio de ejercicio \(K = \EUR{48}\) y vencimiento en 2 años. Hoy quedan \(T-t = 1{,}5\) años hasta el vencimiento.
  
  El valor del contrato forward para el comprador (posición larga) en el momento \(t\) es:
  \[
  F_t = S_t - K(1+r_f)^{-(T-t)}
  \]
  
  Sustituyendo los valores:
  \[
  F_t = 50 - 48(1{,}04)^{-1{,}5} = 50 - 48 \times 0{,}9426 = 50 - 45{,}24 = \EUR{4{,}76}
  \]
  
  \textbf{Interpretación:} El contrato forward tiene un valor positivo de \EUR{4{,}76} para el comprador porque el precio actual de la acción (\EUR{50}) es mayor que el valor presente del precio de ejercicio acordado (\EUR{45{,}24}). Si el comprador vendiera hoy su posición forward, recibiría \EUR{4{,}76}.
  
  \item \textbf{Beneficio neto de la opción call:}
  
  El beneficio de una opción call al vencimiento es:
  \[
  \text{Payoff call} = \max(S_T - E, 0)
  \]
  
  donde \(S_T\) es el precio al vencimiento y \(E = \EUR{48}\) es el precio de ejercicio.
  
  El beneficio neto incluye la prima pagada:
  \[
  \text{Beneficio neto} = \max(S_T - E, 0) - \text{Prima}
  \]
  
  Calculamos para cada escenario:
  \begin{itemize}
    \item \textbf{(i) } \(S_T = \EUR{52}\): 
    \[
    \text{Payoff} = \max(52 - 48, 0) = \EUR{4}
    \]
    \[
    \text{Beneficio neto} = 4 - 3{,}50 = \EUR{0{,}50}
    \]
    
    \item \textbf{(ii) } \(S_T = \EUR{48}\): 
    \[
    \text{Payoff} = \max(48 - 48, 0) = \EUR{0}
    \]
    \[
    \text{Beneficio neto} = 0 - 3{,}50 = -\EUR{3{,}50}
    \]
    
    \item \textbf{(iii) } \(S_T = \EUR{44}\): 
    \[
    \text{Payoff} = \max(44 - 48, 0) = \EUR{0}
    \]
    \[
    \text{Beneficio neto} = 0 - 3{,}50 = -\EUR{3{,}50}
    \]
  \end{itemize}
  
  \textbf{Observación:} La opción call solo genera beneficios cuando el precio al vencimiento supera el precio de ejercicio más la prima (\EUR{51{,}50}). Por debajo de ese nivel, la pérdida máxima está limitada a la prima pagada.
  
  \item \textbf{Beneficio neto total con call y put:}
  
  El beneficio de una opción put al vencimiento es:
  \[
  \text{Payoff put} = \max(E - S_T, 0)
  \]
  
  Con \(S_T = \EUR{44}\) y \(E = \EUR{48}\):
  \[
  \text{Payoff call} = \max(44 - 48, 0) = \EUR{0}
  \]
  \[
  \text{Payoff put} = \max(48 - 44, 0) = \EUR{4}
  \]
  
  El beneficio neto total considerando ambas primas:
  \[
  \text{Beneficio neto total} = 0 + 4 - 3{,}50 - 2{,}00 = -\EUR{1{,}50}
  \]
  
  \textbf{Interpretación:} Aunque la put genera un payoff de \EUR{4}, el inversor tiene una pérdida neta de \EUR{1{,}50} debido a las primas pagadas por ambas opciones (\EUR{5{,}50} en total). Esta estrategia (comprar call y put con el mismo strike) se denomina ``straddle'' y es rentable cuando el precio se mueve significativamente en cualquier dirección, pero requiere movimientos suficientemente grandes para compensar las primas pagadas.
\end{enumerate}
\end{solucion}

 \end{document}
